\subsection{The finite gate with all zero-support places included}

\begin{result}{E.147}{Proved}{An exact finite everywhere-ordinary extension criterion}
Use \(Q_{n-1},Q_n,B_n\) and the Frattini hypothesis
at five and eleven from E.146. Suppose the old set
\(S\) contains \(2,3,5,11,\infty\), with every
finite prime other than three congruent to minus
one modulo three. Fix actual point targets at
the old places other than three. At three fix
an ordinary point \(t_3\), and assume that its
abelianization modulo \(3^{n-1}\) agrees with
the given lower homomorphism.

Put
\[
 M=|Q_n|,\qquad H=(4M+3)^2,\qquad
 T=\{p\notin S:p\text{ prime},\ p\le H\}.
 \tag{E.147.1}
\]
At each \(p\in T\), choose a good-reduction
point target, or allow a choice among such
targets. Only its reduction in
\(\mathbb F_p\setminus\{0,1\}\) matters.
Then the following are equivalent:
\begin{enumerate}[label=(\roman*)]
\item The fixed lower homomorphism has a lift
to \((Q_n/B_n)\rtimes D\), unramified outside
\(S\), matching the old point targets away
from three and the chosen good-reduction
targets at \(T\).
\item It has a lift to \(Q_n\rtimes D\),
with the same old and \(T\)-point matches,
whose abelianization at three equals that
of \(t_3\) modulo \(3^n\), and whose local
class is an actual point image at every
prime different from three. The lift is
unramified outside \(S\) and at most three
new primes, all greater than \(H\), all
congruent to two modulo nine.
\end{enumerate}
The new primes may avoid any additional
prescribed finite set. Their valuation
patterns can be chosen among
\[
 (m,0),\qquad (0,m),\qquad (-m,-m),
       \qquad m\in\{3^{n-1},2\cdot3^{n-1}\}.
 \tag{E.147.2}
\]
In particular the finite abelianization is
locally the reduction of an ordinary pair
at every place.

If a full point match at three is required
instead, the same equivalence holds after
adding that point condition to (i).
As in E.146, actual integer valuations are
prescribed external data; the finite quotient
detects them only modulo \(3^n\).
No compatible infinite choice of the quotient
lifts, or one completed global ordinary pair,
is asserted.
\end{result}
\begin{proof}
We first prove the uniform bound underlying
the finite set \(T\).
Let \(p\ne3\), and let \(s\) be any unramified
local section in the finite geometric quotient
\(Q_n\). The prime-to-\(p\) specialization
comparison for the smooth tripod identifies
it with a section over \(\mathbb F_p\), retaining
the marked boundary inertia
\citep[XIII, Corollary 2.9, \S2.10,
Lemma 2.11, pp.288--290]{RootSGA1}.
Good-reduction point sections correspond
under the same comparison.

The closed subgroup \(s(G_{\mathbb F_p})\)
in \(Q_n\rtimes G_{\mathbb F_p}\) defines
a degree-\(M\) cover \(V_s\) of the tripod
over \(\mathbb F_p\).
Its geometric cover is connected and Galois
with group \(Q_n\), and it is tamely ramified
only above zero, one, and infinity.
This is also the usual twist of the marked
geometric cover by the section cocycle.
Let \(\overline V_s\) be its smooth projective
compactification, of genus \(g\), and let
\(b\) be the number of geometric points
removed over the three boundary points.
Tame Riemann--Hurwitz gives
\[
          2g-2\le M,\qquad b\le3M.
 \tag{E.147.3}
\]
These inequalities hold for every twist:
both quantities are geometric invariants.
The Weil bound
\citep[(1.5.1)--(1.6), pp.275--276]{RootDeligne1974}
therefore gives
\[
 \#V_s(\mathbb F_p)
       \ge p+1-2g\sqrt p-b
       \ge p+1-(M+2)\sqrt p-3M>0
                 \quad\text{if }p>(4M+3)^2.
 \tag{E.147.4}
\]
A point of \(V_s(\mathbb F_p)\) maps to an
interior point whose decomposition section
is geometrically conjugate to \(s\).
Indeed a lift of a rational point to the
cover is precisely containment of its
decomposition group in a conjugate of
\(s(G_{\mathbb F_p})\), and both groups
project isomorphically to \(G_{\mathbb F_p}\).
Lifting that interior residue point to
\(\mathbb Q_p\) proves that every unramified
class at \(p>H\) is an actual good-reduction
point image in \(Q_n\).
At a fixed \(p\ne3\), good-reduction point
images depend only on their residue point.
Thus the possibilities at \(T\) are finite.

Necessity of (i) now follows from E.144:
quotienting by \(B_n\) kills every new
ordinary inertia and retains all the old
and \(T\)-point conditions. The abelian
condition at three reduces to the already
specified lower abelianization.

Assume (i), and fix a quotient lift
\(\bar\rho\) with all its local aligning
frames. As in E.146, the obstruction to
a \(B_n=\mu_3^2\) lift over the old \(S\)
lies in the corresponding Brauer group.
Its local invariants away from three
vanish because the old point targets lift;
its invariant at three is then zero by
the sum formula. Hence there is a global
lift \(t\) unramified outside \(S\).
If a full point at three was prescribed,
this is the same Brauer argument as E.146.

Use the \(\kappa(q_i)\) to correct the
old-support point discrepancies away
from three. Next compare the abelianization
at three with the desired level-\(3^n\)
abelian class of \(t_3\).
The map from \(B_n\) to the new abelian
digit kernel is an isomorphism of
Galois modules:
\[
 B_n\xrightarrow{\sim}
 \ker\bigl((\mathbb Z/3^n(1))^2
          \longrightarrow(\mathbb Z/3^{n-1}(1))^2\bigr).
 \tag{E.147.5}
\]
The local difference reduces to zero
modulo \(3^{n-1}\), so the coefficient
exact sequence gives a correcting class
in \(H^1(\mathbb Q_3,B_n)\).
There is no ambiguity from local
invariants of \(\mu_{3^{n-1}}\), since
\(\mathbb Q_3\) contains no nontrivial
three-power root of unity.
Use the two \(\kappa(3)\) coefficients
to correct its valuation coordinates.
There remain two unit coordinates
\[
                         d=(d_x,d_y)\in\mathbb F_3^2
 \tag{E.147.6}
\]
relative to the unit basis \(\kappa(2)\).
For the stronger full-point version,
perform this same operation on the
complete \(B_n\)-discrepancy at three,
as in E.146.

At every \(p\in T\), the global lift,
the intervening twists, and the chosen
good-reduction point are unramified.
Their aligned kernel discrepancy is
therefore an unramified Kummer class.
For \(p\equiv-1\pmod3\) this group is zero.
For \(p\equiv1\pmod3\), it is the
one-dimensional unit quotient
\[
 \mathbb Z_p^\times/(\mathbb Z_p^\times)^3
             \simeq\mathbb F_3
 \tag{E.147.7}
\]
in each of the two coordinates.
Write its desired correction as
\(r_p\in\mathbb F_3^2\), after fixing
a unit generator.

Choose at most three nonzero vectors
\(u_j\), each a multiple of one of
\[
             e_x,\quad e_y,\quad -e_x-e_y,
 \tag{E.147.8}
\]
such that they span \(\mathbb F_3^2\)
and sum to \(d\).
Here is an explicit verification.
If \(d_x,d_y\) are distinct and both
nonzero, take \(d_xe_x,d_ye_y\).
Otherwise choose \(c\in\{1,2\}\) with
\(c+d_x\ne0\) and \(c+d_y\ne0\), and take
\[
 (c+d_x)e_x,\quad(c+d_y)e_y,\quad
                            -c(e_x+e_y).
 \tag{E.147.9}
\]
The excluded case is exactly the case
where the two requirements rule out both
possible \(c\). These vectors have the
claimed sum, and the first two span.
This includes \(d=0\), when all three
peripheral directions with coefficient
one have sum zero.
For each \(p\in T\) with \(p\equiv1\pmod3\),
choose scalars \(\beta_{j,p}\) solving
\[
                      \sum_j u_j\beta_{j,p}=r_p;
 \tag{E.147.10}
\]
spanning guarantees their existence.

We show that these unit prescriptions
are compatible with every required
cusp Frobenius.
Enlarge the actual arithmetic data to
retain \(\mu_9\) and all three finite
tangential cusp paths. The infinity
path satisfies the same unramifiedness
condition as the other two: use the
coordinate \(1/t\) and the unit tangent,
or the tripod automorphism interchanging
zero and infinity in the actual path
statement of E.128.
All these data are unramified outside
three.
The five and eleven inertia vectors
still give the full geometric image
over this enlarged arithmetic quotient.

Let \(L\) be this finite joint quotient
field of \(\bar\rho\), the arithmetic
action, and the cusp paths.
It is unramified at every prime of \(T\).
For each \(p\in T\), \(p\equiv1\pmod3\),
let \(C_p\) be the cyclic cubic subfield
of \(\mathbb Q(\zeta_p)\).
Its Frobenius at a prime \(q\ne p\)
records exactly the unit cube class
of \(q\) in (E.147.7), after choosing
the corresponding generator.
Put \(C_T=\prod_{p\in T,\ p\equiv1(3)}C_p\).
Then
\[
                          L\cap C_T=\mathbb Q .
 \tag{E.147.11}
\]
Indeed \(C_T/\mathbb Q\) has group the
product of its independent cubic inertia
groups at the distinct primes \(p\).
Any subextension unramified at all these
primes is therefore trivial.
The intersection is unramified there,
because \(L\) is, proving (E.147.11).

For each vector \(u_j\), use the full
geometric image to choose the
corresponding cusp Frobenius in \(L\)
above an arithmetic element with
cyclotomic value two modulo nine.
Independently prescribe the Frobenius
in \(C_p\) to be \(\beta_{j,p}\) for
every \(p\in T\).
Equation (E.147.11) and Chebotarev give
distinct primes \(q_j>H\), outside every
additional finite forbidden set, with
all these conditions simultaneously.
In particular \(q_j\equiv2\pmod9\).
Twist \(t\) by the global kernel cocycle
\[
                         \sum_j u_j\kappa(q_j).
 \tag{E.147.12}
\]
Its contribution at three is
\(\sum_j u_j\kappa(2)=d\kappa(2)\).
Its contribution at each \(p\in T\)
is exactly (E.147.10).
At each old support prime different
from three it is zero, because all
the new \(q_j\) are local units and
all such units are cubes there.
Consequently every old and \(T\)
point condition is attained, as well
as the required condition at three.

At a new prime \(q_j\), the same
calculation as E.146 applies:
the quotient Frobenius is the selected
cusp, \(H^1_{\rm nr}(\mathbb Q_{q_j},B_n)=0\),
and the kernel class is determined by
its specified peripheral inertia.
For directions \(e_x,e_y\) this is the
actual point \(q_j^m\) or \(1-q_j^m\).
For direction \(-e_x-e_y\), E.144's
infinity-power identity identifies
the kernel generator with the marked
infinity peripheral power.
The actual point may be taken as
\(q_j^{-m}\); its two valuations are
\((-m,-m)\).
The local cusp twist and A.31 again
identify the entire local class,
including Frobenius.
All other new-prime twists are locally
unit cubes and do not alter that class.

The final lift is unramified outside
\(S\cup\{q_j\}\).
At every remaining prime outside \(S\)
and \(T\), the bound (E.147.4) applies,
so its local class is a good-reduction
point image. At real infinity the
old point condition is unchanged,
since odd-order kernel first cohomology
vanishes there.
This proves (ii) and the equivalence.
The construction leaves exactly the
quotient matching test (i); it does not
provide an induction showing that this
test can always be passed.
\end{proof}
